% File configurations
\documentclass[12pt, a4paper]{paper}
\usepackage{babel}
\usepackage{graphicx}
\usepackage[margin=2.5cm]{geometry}
\usepackage{listings}
\usepackage{multicol}
\title{Javascript \\
  \hfill\includegraphics[height=3cm]{../images/universidade.png}
  \vspace{-3cm}
}
\subtitle{Do básico ao avançado seguindo a MDN web docs}
\author{João Vitor Rezende Moura - 1221184773}
\institution{UNIT - Universidade Tiradentes}

% Specific elements created


\begin{document}
% Document specific elements
\maketitle

\hrule
\begin{abstract}
  abstract
\end{abstract}
\vspace{-0.6cm} % vspace 
\begin{keywords}
  Javascript, Front-End, UI, media-queries, web
\end{keywords}
\hrule

% ------------------ Content -------------------------------
\section{Primeiros passos com javascript} % (fold)
\label{sec:Primeiros passos com javascript}

% section Primeiros passos com javascript (end)

\section{Criando elementos em javascript} % (fold)
\label{sec:Criando elementos em javascript}

\subsection{Condicionais dentro do javascript} % (fold)
\label{sub:Condicionais dentro do javascript}
Dentro de qualquer linguagem de programação, temos que tomar decisões, 
sejam elas baseadas nas informações providas pelo usuário, nas informações 
que conseguimos captar dele, ou seja por algum outro motivo. Esse tipo de 
escolha, pode ser feito por meio de uma estrutura que existem em todas as 
linguagens de programação, e que recebe o nome de \textbf{\textit{condicionais}}.\par

Elas permitem que a partir de determinada condição que deve ser analisada,
executar-
mos uma determinada ação, ou outra. Um exemplo de código javascript 
que possui essa determinada estrutura, é a seguinte:

\begin{verbatim}
  if(condição){
    // codigo executado se a condição for verdadeira
  }else{
    // código executado se a condição for falsa
  }
\end{verbatim}

A priori, temos que as palavras \texttt{if} e \texttt{else} são reservadas dentro do 
javascript para a construção de estruturas condicionais, portanto
temos que dentro dos parênteses, temos a condição que deve ser analisada, 
e dentro das chaves, os blocos de códigos a serem executados condicionalmente.

A posteriori, temos que as vezes as comparações condicionais de dois fatores
podem não ser suficientes para nós conseguirmos analisar o código e executá-lo 
da maneira desejada, portanto, temos a necessidade de comparar um determinado valor 
a várias possibilidades as quais ele pode assumir, e assim surge o \texttt{else if}:

\begin{verbatim}
  if(primeiraCondicao){
    // executar código do primeiro bloco
  }else if(segundaCondicao){
    // executar o segundo bloco de código
  }else{
    // executar bloco de código quando condições não forem atingidas
  }
\end{verbatim}

\subsection{Aninhando condicionais, e estabelecendo condições} % (fold)
\label{sub:Aninhando condicionais, e estabelecendo condições}
Temos que as condições presentes analisandas dentro das condicionais,
devem assumir valores booleanos, ou seja, de verdadeiro ou falso. Entretanto,
podemos criar análises numéricas, decimais, e de outros tipos primitivos, que 
possuem como resultado, um valor booleano, e inserir como a condição que controla
o fluxo da condicional, como no seguinte caso:

\begin{verbatim}
  let condicao = 15 < 20

  if(condicao){
    // executar código se a condicao for verdadeira
  }else{
    // executar código se a condicao for falsa

\end{verbatim}

Temos no exemplo anterior, que uma comparação lógica entre dois elementos numéricos
gerou um valor booleano, e isso pode ser feito de várias maneiras por meio dos operadores lógicos, ou por meio de operadores de comparação. Os operadores que possibilitam isso, são os de comparação matemática, e operadores booleanos, que são o \texttt{\&\&} e o \texttt{||}, que são de disjunção e conjunção respectivamente

Podemos também aninhar condicionais, ou seja, dentro de alguma determinada condi-
cional, fazer outra condicional que pode comparar o mesmo valor, ou um valor que não possua
conexão com o primeiro comparado.
% subsection Aninhando condicionais, e estabelecendo condições (end)

\subsection{Instrução switch} % (fold)
\label{sub:Instrução switch}
As isntruções condicionais são muito úteis, entretanto, quando temos uma variável
que pode assumir diversos valores, isso cria uma cadenia de condicionais que é dificil
de manter no código, e difícil de entender, portanto, é recomendando que quando 
comparamos condicçoes que podem assumir diversos valores, que usemos 
a estrutura \texttt{switch}.

\begin{verbatim}
  switch(expression){
    case choice1:
      // run code and break the comparation
      break;
    case choice2:
      // run code and break the comparation
      break;
    default:
      // only need to run the code actually
  } 
\end{verbatim}

Nesse caso, temos que dois casos existem, e que quebramos a comparação como última 
instrução do caso o qual o mesmo corresponde. Isso se deve ao fato de que o switch case 
pode desejar que uma expressão, encaixe em mais de um caso, sendo assim, o mesmo deve 
passar pelos dois casos, e executar os dois blocos de isntrução, e isso é feito não 
quebrando a comparação dos casos, ou seja, não adicionando a instrução \texttt{break}.
Caso nenhum dos casos seja verdadeiro, ou caso o que seja, não possua a instrução de 
quebrar a comparação, e seja o único verdadeiro, a expressão vai cair no caso padrão, 
e após executar as instruções, vai quebrar a comparação por ser o último caso 
propositadamente.
% subsection Instrução switch (end)

\subsection{Operador ternário} % (fold)
\label{sub:Operador ternário}
Temos que as comparações são elementos essenciais dentro da programação, entretanto 
elas ocupam muito espaço visual, e algumas comparações são tão pequenas que são quase
"banais", pra isso existe o operador ternário, que permite uma comparação de um 
elemento, que pode receber um de dois valores se sua condição for verdadeira, e essa 
instrução é feita toda em um única linha.

\begin{verbatim}
  (condição) ? se a condição for verdadeira : se a condição for falsa;
\end{verbatim}
% subsection Operador ternário (end)

\section{Loops} % (fold)
\label{sec:Loops}
Os loops. são estruturas que permitem repetição de determinadas partes do código, 
até que uma determinada condição que pode ser pré-definida, ou mudar conforme a 
execução do código, seja realizada. Eles possuem uma estrutura padrão, mas podem se 
modificar conforme as necessidades e peculiaridades de cada código. Por padrão 
os laços dentro do Javascript são feitos da seguinte forma: 

\begin{verbatim}
for(inicializador; condição de saída; expressão final){
  //código
}

\end{verbatim}

\subsection{Saindo dos laços por meio da instrução break}
\label{sub:Saindo dos laços por meio da instrução break} % (fold)
Apesar dos laços possuirem uma quantidade definida de vezes as quais eles repetem um 
conjunto específico de ações, eles podem sofrer alterações no seu comportamente e 
podem também sofrerem desvios por meio de condicionais que podem ser colocadas dentro 
dos laços, com intuito de caso uma certa condição seja apresentada, aquela interação 
e as seguintes sejam canceladas.

\begin{verbatim}
for(inicializador; condição de saida; expressao final){

  // codigo
  if(condicao){
    break;
  }
}

\end{verbatim}


\textbf{Obs:} O mesmo pode ser feito com a instrução \textbf{continue;}, so que 
diferentemente, ele não vai cacelar todas as interações, ele só vai cancelar a atual, 
e prosseguir com as seguintes.


\subsection{While e do...while}
\label{While e do...while}
Essas instruções são parecidas com os laços, entretanto, a diferença é que elas 
apresentam em si, o fato de que ela não necessariamente possuem dentro de si uma 
condição de saída, mas sim uma condição de continuidade, em que caso ela seja 
satisfeita, aquele laço será repetido até que essa mesma condição seja quebrada. 

Temos que essa instrução pode ser apresentada e executada de duas maneiras, na primeira,
ele confere a condicional, e se for verdadeira, ele executa um bloco de código e parte 
para a próxima verificação. Por outro lado, podemos fazer com que ele primeiro execute 
a ação, e somente posteriormente verifique as condicionais, como mostrados os exemplos
seguintes: 

\begin{verbatim}
while(condicao){

  // executar código

}

do{
  \\ executar código

}while(condição)
\end{verbatim}
% section Loops (end)


\section{Funções - Blocos reutilizáveis de código} % (fold)
\label{sec:Funções - Blocos reutilizáveis de código}
Funções são rotinas de código pré-definidas que podem ser utilizadas diversas vezes, e 
que são agrupadas dentro de uma mesma chamada, ou seja, invocando somente um comando 
dentro do código, todas aquelas ações correspondentes presentes dentro da nossa 
função, será executadas.

Elas podem vir de natureza da própria linguagem, serem chamadas a partir de objetos, ou 
serem criadas pelo desenvolvedor para executar uma rotina específica de funções que 
existem, elas podem aparecer das seguintes formas: 

\begin{verbatim}
  // embutida do navegador    
  console.log();
  // por meio de objetos 
  MyArray.join("");
  // funções criadas pelo desenvolvedor 
  function draw(){
    console.log("desenhando");
  }
  draw();
\end{verbatim}

\subsection{Funções vs. Métodos} % (fold)
\label{sub:Funções vs. Métodos}
Uma coisa técnicas mas que pode ser pontuada, é o fato de que as funções embutidas 
nos navegadores na verdade são chamadas de métodos, pois o navegador é um objeto, e 
toda função que referencia um objeto é chamada de método, e as funções não 
necessariamente referenciam objetos. 
% subsection Funções vs. Métodos (end)

As funções dentro do javascript, podem possuir uma característica muito especial, que 
é serem anônimas, ou seja, não possuírem um nome definido que as permitam serem 
invocadas. Normalmente fazemos isso quando criamos uma função junto com um 
manipulador de eventos.

\subsection{Parâmetros de funções} % (fold)
\label{sub:Parâmetros de funções}
Algumas funções requerem parâmetros a serem especificados quando estamos invocando eles,
 esses valores precisam ser inclusos dentro dos parênteses da função, o que é  
 necessário para fazer seu trabalho apropriado. 
% subsection Parâmetros de funções (end)

\subsection{Função escopo e conflitos} % (fold)
\label{sub:Função escopo e conflitos}
Um conceito muito importante quando lidamos com funções é o escopo do seu uso. As 
funções, e suas variáveis possui um escopo, ou seja eles funcionam apenas dentro do 
seu próprio funcionamento, sendo inacessível para outras funções ou de código do lado 
de fora das funções. 

Essa "passagem" de valores só pode ser feito por meio da passagem dos valores do 
escopo local, por meio ad chamada da outra função, e a enviando-a como parâmetro.
 % subsection Função escopo e conflitos (end)

\subsection{Retornando valores} % (fold)
\label{sub:Retornando valores}
As funções elas podem ser entendidas, como rotinas de ações que executam e modificam
dados, portanto, pode-se desejar obter de volta os dados modificados, ou obter os 
dados que foram criados. Para isso, podemos usar a palavra reservada \textbf{return} 
para indicar um dado que deve ser retornado da chamada da função.
% subsection Retornando valores (end)
% section Funções - Blocos reutilizáveis de código (end)

\section{Introdução a eventos.} % (fold)
\label{sec:Introdução a eventos.}
Eventos são ações ou ocorrências que acontecem no sistema que estamos desenvolvendo, no 
qual este te alerta sobre ações para que as quais devemos poder responder de alguma 
determinada forma. Eventos web são disparados dentro da janela do navegador e tendem a 
estarem anexados a lagum item específico nele. 

Os \textbf{Manipuladores de eventos} são blocos de código que são executados pelo 
evento quando ele é disparado, esse elemento também pode ser chamado de 
\textit{event handler}. Um exemplo de seu uso pode ser visto a seguir:

\begin{verbatim}
HTML 
<button>Change color</button>


JS
var button = document.querySelector("button")

function random(number){
  return Math.floor(Math.random() * (number + 1));
}

button.onClick = function(){
  var randomColor = "rgb(" + random(255) + 
                       "," + random(255) + "," + random(255) + ")"
  document.body.style.backgroundColor = randomColor;


\end{verbatim}


\subsection{Não são apenas páginas web} % (fold)
\label{sub:Não são apenas páginas web}
O modelo de eventos no JavaScript para web pages difere dos outros modelos de eventos 
próprios que as outras linguagens de programação podem conter. O \textbf{Node.js} é 
um interpretador de código JavaScript que permite o desenvolvimento de aplicações do 
lado do servidor. O event model do Node.js depende dos ouvintes para escutar eventos e 
emissores para emitir eventos periodicamente.

\subsection{Maneiras de usar os eventos Web} % (fold)
\label{sub:Maneiras de usar os eventos Web}
Hpa várias maneiras de adicionar código com event listeners a uma determinada página 
web, mas a forma mais comum disso ser feita é por meio de propriedades nas páginas as 
quais são adicionadas os manipuladores de eventos. As \textbf{propriedades} podem ser 
usadas para conter código que manipulam os eventos. Essas propriedades são 
essencialmente iguais a qualquer outra disponível, mas com um tipo especial o qual 
pertencem.
% subsection Maneiras de usar os eventos Web (end)

\subsection{addEventListener e removeEventListener} % (fold)
\label{sub:addEventListener e removeEventListener}
Temos que os mecanismos de eventos são definidos nas especificações de eventos que 
ficam presentes dentro da DOM(Document Object Model), o qual fornece aos navegadores 
uma nova função, que é a de adição de eventlisteners. Dentro dessas funções, 
especificamos o nome do evento para o qual queremos registrar esse manipulador, e o 
código que compreende a função do manipulador que queremos executar em resposta a ele. 

\begin{verbatim}
let button = document.querySelector("button");

button.addEventListener("click", => (){
  document.body.style.backgroundColor = randomcolor...;
});


// Se quisermos remover o event listener, podemos usar o seguitne código 

button.removeEventListener("click", colorChange)
\end{verbatim}

Uma observação que deve ser feita sore o event listener e sua remoção, é que para a 
mesma seja feita, precisamos especificar o nome da função responsável pelo event 
listener.
% subsection addEventListener e removeEventListener (end)

\subsection{Objetos de Eventos} % (fold)
\label{sub:Objetos de Eventos}
Dentro de uma função a qual manipulaos o evento podemos desejar verificar um parâmetro 
chamado \texttt{event}. É um objeto representante do evento em si, e é passado 
automaticamente para os manipuladores de eventos para fornecer recursos e informações 
extras.
Esse evento possui o atributo target, que é sempre uma referência ao elemento em que o 
evento acabou de ocorrer. Ele é útil quando desejamos definir o mesmo manipulador de 
eventos em vários elementos e fazer algo com todos eles quando ocorre um evento neles. 
Isso pode ser feito com o método \texttt{querySelectorAll()}.

Outra coia importante é que podemos implementar uma verificação simples dentro do 
manipulador de eventos, como dentro do \texttt{onsubmit}, por exemplo, que verifica 
se a função no objeto do evento tem campos vazios.

\begin{verbatim}
var form = document.querySelector("form");
var fname = document.getElementById("fname");
var lname = document.getElementById("lname");
var submit = document.getElementById("submit");
var para = document.querySelector("p");

form.onsubmit = function (e) {
  if (fname.value === "" || lname.value === "") {
    e.preventDefault();
    para.textContent = "You need to fill in both names!";
  }
};
\end{verbatim}
% subsection Objetos de Eventos (end)

\subsection{Borbulhamento e Captura de eventos} % (fold)
\label{sub:Borbulhamento e Captura de eventos}
Borbulhamento e captura de eventos são dois mecanismos que descrevem oque acontece 
quando dos manipuladores do mesmo tipo de evento são ativados em um elemento. Quando 
temos dois manipuladores de eventos, quando um evento é acionado em um elemento que 
possui elementos pai, os navegadores executam duas fases diferentes, que são as fases 
de captura de bubbling. 

\subsubsection{Fase de Captura} % (fold)
\label{sec:Fase de Captura}
Na fase de captura, o nabegador verifica se o ancestral mais externo do elemento tem 
um manipulador de eventos registrado nele na fase de captura e executa em caso 
afirmativo, o manipulador verificado é o \texttt{onclick}. Em seguida ele passa para 
o próximo elemento dentro do html e faz a mesma coisa, e assim por diante até alcançar o elemento que foi realmente clicado.
% subsubsection Fase de Captura (end)

\subsubsection{Fase de bubbling} % (fold)
\label{sec:Fase de bubbling}
O navegador verifica se o elemento que foi clicado realmente tem o manipulador de 
eventos específico registrado nele e o executa em caso afirmativo. Em seguida ele passa 
para o próximo elemento até alcançar o elemento html
% subsubsection Fase de bubbling (end)

Esse tipo de problema pode ser corrigido com o objeto padrão chamado o método 
\texttt{stopPropagation()}, que quando invocado no objeto de evento de um manipulador 
faz com que o manipulador seja executado, mas o evento não borbulha mais acima da 
cadeia.
% subsection Borbulhamento e Captura de eventos (end)
% section Introdução a eventos. (end)

\section{Intodução a Objetos em JavaScript} % (fold)
\label{sec:Intodução a Objetos em JavaScript}
\subsection{Noções básicas de Objetos} % (fold)
\label{sub:Noções básicas de Objetos}
Um objeto é uma coleção de dados e/ou funcionalidades relacionadas. Eles podem ser 
armazenados dentro de variáveis, como valores descritos, para criar um objeto vazio, 
podemos fazer da seguinte maneira: 

\begin{verbatim}
  var pessoa = {};

  // Console
  [object Object]
\end{verbatim}

Entretanto, podemos criar objetos pré-definidos para o uso, possuindo atributos e 
métodos já instanciados dentro da sua criação, como mostrado a seguir: 

\begin{verbatim}
  var pessoa = {
    nome: ["",""],
    idade: 32,
    sexo: "masculino",
    ...
  }
\end{verbatim}

para acessarmos as propriedades e métodos de algum determinado objeto, devemos usar 
a notação de ponto, que é uma notação que permite acessarmos essas informações dentro 
do nosso objeto, como mostrado a seguir:

\begin{verbatim}
  pessoa.nome;
  pessoa.sexo;
\end{verbatim}

Ooutra forma de notação existente e que podemos usar também é a notação de colchetes,
onde especificamos que atributo específico do nosso objeto queremos modificar ou 
acessar dizendo seu nome dentro dos colchentes, como mostrado a seguir:

\begin{verbatim}
  pessoa["nome"];
  pessoa["sexo"];
\end{verbatim}

A vantagem da notação de colchete comparada a notação de ponto, é o fato de que a 
notação de colchete permite a criação de um novo atrbituo e sua instanciação ao mesmo 
tempo, sem a necessidade de recriar o objeto.
% subsection Noções básicas de Objetos (end)

\subsection{Escopo de Classe e de Instância} % (fold)
\label{sub:Escopo de Classe e de Instância}
Quando criamos um objeto dentro do jvascript, podemos ter variáveis as quais devem 
pertencer a toda a classe, independente da instância, e podemos ter elementos os 
quais necessitamos que eles possam possuir múltiplos valores, e que possam ser 
acessados de forma individual, que são as variáveis de instância, elas podem ser 
acessadas por meio da palavra chave \texttt{this}. 

\begin{verbatim}
  var pessoa = {
    nome: "Chris";
    saudacao: function(){
      alert("Oi, meu nome é: " + this.nome + "!")
    }
  }
\end{verbatim}
% subsection Escopo de Classe e de Instância (end)

\subsection{Protótipo de Objetos } % (fold)
\label{sub:Protótipo de Objetos }

% subsection Protótipo de Objetos  (end)



% section Intodução a Objetos em JavaScript (end)



\end{document}




